\section{统一视角：反馈从哪里来}

自我改进系统需要一个信号告诉模型“刚才做得怎样”。本讲比较三类反馈：

\begin{enumerate}
    \item \textbf{工具观察（tool observation）}：搜索结果、网页内容或环境状态；
    \item \textbf{执行反馈（execution feedback）}：编译错误、单元测试、运行时行为；
    \item \textbf{AI feedback}：模型依据人类编写的原则批评、修订并比较回答。
\end{enumerate}

\videofigure{figures/thought-action-feedback.jpg}{思考指导行动，行动产生反馈，反馈再修正思考。}{00:01:32--00:03:32}

三类方法分别对应 ReAct、RLEF 和 Constitutional AI。它们共享同一闭环：

$$
\text{state}\rightarrow\text{reason / act}\rightarrow\text{feedback}
\rightarrow\text{revision or update}.
$$

\begin{importantbox}{反馈可以只改当前轨迹，也可以改模型参数}
ReAct 主要在推理时读取 observation 并调整下一步；RLEF 把执行结果作为 RL reward 更新代码模型；Constitutional AI 先把 critique--revision 数据用于 SFT，再把 AI preference 用于 RL。区分这两层能避免把“会重试”误认为“已经学会”。
\end{importantbox}

\subsection{评价反馈质量的四个维度}

\begin{itemize}
    \item \textbf{真实性}：是否来自真实环境而非模型自说自话；
    \item \textbf{粒度}：只告诉最终成败，还是指出具体错误；
    \item \textbf{成本}：是否需要人类专家、沙箱执行或额外大模型；
    \item \textbf{可扩展性}：能否自动产生大量训练样本。
\end{itemize}

\subsection{本章小结}

反馈是模型与目标之间的接口。工具和代码提供外部可观测信号，AI feedback 提供可扩展但更主观的规范信号；系统的可靠性取决于这些信号是否真实对应最终目标。

